\section{Introduction}

Data curation for foundation model training combines quality filtering and diversity sampling, but optimal ordering remains an open question. When filtering 70\% of data by quality (V-Information) before sampling 49\% for diversity (Feature Activation Coverage), or applying these steps in reverse, which strategy wins? Conventional wisdom from prior qualitative observations suggests quality-first (QD) dominates at small scale while diversity-first (DQ) should take over at large scale as quality saturates---but our experiments with GPT-2 training across 10K--10M tokens reveal quality-first wins at ALL scales.

This finding matters because practitioners deploying multi-stage curation pipelines need prescriptive guidance on ordering, not just quality and diversity metric selection. Training on 10M tokens with DQ ordering (diversity-first) yields models 1.8 percentage points worse than QD despite intuition that diversity should dominate at large scale. Incorrect ordering wastes training compute and produces suboptimal models even when final dataset size remains constant.

The core challenge is compositional interaction: sequential curation strategies create non-additive effects that prior work on domain mixing assumes to be order-invariant. Quality metrics show diminishing returns as dataset scale increases---DATAMASK (ByteDance 2025) qualitatively observed this saturation---while diversity metrics appear to persist. Yet existing approaches like Data Mixing Laws and RegMix assume data sources mix independently with additive contributions when combining domains, not addressing sequential filtering dependencies. No prior work quantifies quality saturation and diversity persistence trajectories as continuous functions of scale, nor systematically tests whether ordering (QD vs DQ) produces statistically different outcomes.

Our key insight is that \textbf{quality-first filtering appears to act as a diversity-preserving prior} by removing low-value noise before diversity selection, while diversity-first sampling may amplify uninformative variation when applied to unfiltered noisy corpora. Think of quality filtering as removing static from a radio signal before applying equalization (diversity sampling). If you equalize first, you amplify both signal and noise. The quality-first approach maximizes signal-to-noise before pattern coverage optimization.

Building on this insight, we make the following contributions:

\begin{enumerate}
\item \textbf{Quality-gated diversity principle}: We show evidence that diversity sampling is primarily effective on quality-filtered subsets. Quality-first (QD) outperforms diversity-first (DQ) universally across tested scales by +1.5pp to +5.0pp (Cohen's $d=0.76$--$2.48$), even when diversity coefficient exceeds quality coefficient at large scale.

\item \textbf{Saturation and persistence trajectories quantified}: Quality-only improvement decreases from +11pp (10K tokens) to +3.3pp (10M tokens) with slope $-2.62$pp/log-scale ($p=0.008$), providing statistical rigor to DATAMASK's qualitative observations. Diversity persistence is modeled based on FAC literature (simulated trajectory pending full implementation validation).

\item \textbf{Evidence for compositional ordering effects}: Our systematic GPT-2 training experiments across scales suggest that sequential ordering (QD $\neq$ DQ) produces statistically significant compositional interactions in data curation, extending Data Mixing Laws' domain proportion framework to sequential filtering steps.

\item \textbf{Prescriptive guidance for practitioners}: Our findings suggest actionable recommendations for 100M--1B parameter models trained on web corpora at 10K--10M scale---apply quality filtering before diversity sampling to improve model performance while saving compute on DQ experiments.
\end{enumerate}

Our work extends the understanding of how data sources combine by showing that sequential filtering order matters, not just domain proportions. We organize the paper as follows: Section~2 discusses related work in data curation metrics and mixing strategies; Section~3 describes our experimental methodology; Section~4 presents results on saturation trajectories and ordering effects; Section~5 discusses implications and limitations; Section~6 concludes with future directions.
